A trace formula is a way of bringing two descriptions of the same mathematical quantity into contact. In its most elementary form, it says that the sum of the diagonal entries of a matrix is the sum of its eigenvalues, counted with multiplicity. One description uses the entries of an operator; the other uses its action on spectral data. When an operator is built from geometry or arithmetic, equating these descriptions can turn information about one into information about the other.

This chapter explains the pattern before introducing modular forms. We will compute traces of finite matrices, replace matrix entries by a finite-rank kernel, and examine a finite Fourier identity whose continuous analogue is Poisson summation. Those models let us distinguish several questions that will otherwise look deceptively similar: counting eigenvalues, averaging Fourier coefficients, taking Hecke traces, pairing periods, and measuring where eigenfunctions place their mass.

The finite models are mathematical results with proofs. The later automorphic formulas are introduced as destinations: their definitions, constants, convergence conditions, and proofs belong to the chapters that develop them. In particular, an infinite sum of eigenvalues does not become meaningful merely because its finite-dimensional counterpart is easy to calculate.

\subsection*{How to read this chapter}

\hyperref[v1a-ch00-s1]{Section 0.1} uses elementary linear algebra, determinants, and the fundamental theorem of algebra. \hyperref[v1a-ch00-s2]{Section 0.2} introduces a small amount of integration and inner-product notation, with all operator manipulations reduced to finite sums. \hyperref[v1a-ch00-s3]{Section 0.3} proves the finite Fourier model and states the analytic Poisson formula as a preview. \hyperref[v1a-ch00-s4]{Section 0.4} describes the questions answered by the major formulas in this volume. \hyperref[v1a-ch00-s5]{Section 0.5} returns to the finite kernel and shows exactly how ordinary, relative, and weighted constructions differ.

You can read the conceptual passages first and return to the calculations. The calculations are also useful as normalization checks: every more elaborate theory should agree with its finite model in the cases where both apply. Exercises include solutions so that none of the subsequent discussion depends on an argument hidden in an unsolved problem.

By the end, the aim is to recognize what is being summed or integrated, which data are retained by that operation, and what still has to be justified before an infinite-dimensional version can be used. We will then be ready to construct the geometric setting of the book: the upper half-plane and the modular group.

\textit{Editorial status: complete chapter draft for Steven's review; formalization and human correspondence review remain pending.}
